\documentclass[handout]{beamer}
\mode<presentation>
\usetheme[secheader]{Boadilla}
\usecolortheme{whale}

\usepackage[utf8]{inputenc}
\usepackage[T1]{fontenc}
\usepackage[indonesian]{babel}
\usepackage{lmodern}
\usepackage{graphicx}
\usepackage{bookmark}

\setbeamertemplate{footline}
{
	\leavevmode%
	\hbox{%
		\begin{beamercolorbox}[wd=.333333\paperwidth,ht=2.25ex,dp=1ex,center]{author in head/foot}%
			\usebeamerfont{author in head/foot}\insertshortauthor%
		\end{beamercolorbox}%
		\begin{beamercolorbox}[wd=.433333\paperwidth,ht=2.25ex,dp=1ex,center]{title in head/foot}%
			\usebeamerfont{title in head/foot}\insertshorttitle
		\end{beamercolorbox}%
		\begin{beamercolorbox}[wd=.233333\paperwidth,ht=2.25ex,dp=1ex,right]{date in head/foot}%
			\usebeamerfont{date in head/foot}\insertshortdate{}\hspace*{2em} \insertframenumber{} / \inserttotalframenumber\hspace*{2ex}
	\end{beamercolorbox}}%
	\vskip0pt%
}

\title[Bab 6 -- Form \& Media]{Pemrograman Web}
\subtitle{Bab 6: HTML Form, Media, Validasi Dasar, dan Aksesibilitas Form}
\author{Cecep Suwanda, S.Si., M.Kom.}
\institute{Universitas Bale Bandung}
\date{\today}

\begin{document}

% Slide Judul
\begin{frame}
	\titlepage
\end{frame}

% Outline dan CPMK
\begin{frame}{Outline Bab 6}
	\begin{block}{Sub-CPMK 2.1}
		Membuat halaman web dengan HTML5 (struktur, elemen semantik, form, media).
	\end{block}
	\begin{block}{Sub-CPMK 1.3}
		Menerapkan prinsip aksesibilitas (a11y) dan pengalaman pengguna (UX) dalam desain antarmuka.
	\end{block}
	\begin{itemize}
		\item Tabel HTML: Dasar (Baris, Sel, Header)
		\item Tabel: Pengelompokan, Scope, dan Aksesibilitas
		\item Form HTML: Elemen form, input, label, dan button
		\item Form: Tipe input, select, textarea, dan validasi dasar
		\item Struktur Halaman dengan Elemen Semantik
		\item Media: Audio dan Video
		\item Gambar Vektor (SVG) dalam HTML
		\item Embedding: iframe, object, dan embed
		\item Aksesibilitas (a11y): Prinsip dan Praktik
		\item Aksesibilitas Tabel dan Form
		\item Validasi HTML dan Alat Validator
		\item Debugging HTML: Menemukan dan Memperbaiki Error
		\item Best Practices HTML dan Standar (WHATWG / W3C)
	\end{itemize}
\end{frame}

% ========== SECTION 6-1: Tabel Dasar ==========
\begin{frame}{Tabel HTML: Dasar (Baris, Sel, Header)}
	\begin{itemize}
		\item Elemen \texttt{<table>}: wadah tabel untuk data terstruktur (baris-kolom).
		\item \texttt{<tr>}: baris tabel; \texttt{<td>}: sel data; \texttt{<th>}: sel header (tebal, semantik).
		\item \texttt{<caption>}: judul tabel untuk aksesibilitas.
		\item \texttt{colspan} dan \texttt{rowspan}: menggabungkan sel horizontal/vertikal.
	\end{itemize}
	\begin{block}{Kapan tabel?}
		Data tabular (jadwal, daftar nilai, statistik). Bukan untuk layout halaman---gunakan CSS Flexbox/Grid.
	\end{block}
\end{frame}

% ========== SECTION 6-2: Tabel Pengelompokan ==========
\begin{frame}{Tabel: Pengelompokan, Scope, dan Aksesibilitas}
	\begin{itemize}
		\item \texttt{<thead>}: baris header; \texttt{<tbody>}: baris data utama; \texttt{<tfoot>}: ringkasan/total.
		\item \texttt{scope="col"}: header untuk kolom; \texttt{scope="row"}: header untuk baris.
		\item \texttt{id} + \texttt{headers}: hubungan eksplisit sel--header (tabel kompleks).
		\item \texttt{<caption>}: pembaca layar bacakan caption dahulu; singkat dan informatif.
	\end{itemize}
	\begin{block}{Hindari}
		Tabel bersarang; tabel lebar di mobile tanpa alternatif.
	\end{block}
\end{frame}

% ========== SECTION 6-3: Form Dasar ==========
\begin{frame}{Form HTML: Elemen form, input, label, dan button}
	\begin{itemize}
		\item \texttt{<form>}: wadah; \texttt{action} (URL tujuan), \texttt{method} (GET/POST).
		\item GET: data via URL query (pencarian); POST: data di body request (data sensitif).
		\item \texttt{<label>}: pasangan input; \texttt{for} pada label = \texttt{id} pada input.
		\item \texttt{<button>}: \texttt{type="submit"} (kirim), \texttt{type="reset"} (kosongkan), \texttt{type="button"} (JS).
	\end{itemize}
	\begin{block}{Label}
		Klik label $\to$ fokus/aktifkan input; penting untuk a11y dan UX (area klik lebih besar).
	\end{block}
\end{frame}

% ========== SECTION 6-4: Tipe Input dan Validasi ==========
\begin{frame}{Form: Tipe input, select, textarea, dan validasi}
	\begin{itemize}
		\item Tipe HTML5: \texttt{text}, \texttt{email}, \texttt{password}, \texttt{number}, \texttt{date}, \texttt{range}, \texttt{checkbox}, \texttt{radio}, \texttt{file}.
		\item \texttt{<select>}, \texttt{<option>}, \texttt{<optgroup>}; \texttt{<textarea>} untuk teks panjang.
		\item Validasi: \texttt{required}, \texttt{minlength}/\texttt{maxlength}, \texttt{pattern}, \texttt{min}/\texttt{max}.
		\item \texttt{<fieldset>} dan \texttt{<legend>}: pengelompokan field; penting untuk radio/checkbox.
	\end{itemize}
\end{frame}

\begin{frame}{Validasi Client-side vs Server-side}
	\begin{block}{Client-side (HTML/JS)}
		Umpan balik instan; mengurangi request gagal; \textbf{dapat dilewati} oleh pengguna.
	\end{block}
	\begin{block}{Server-side (PHP/dll.)}
		Menjamin keamanan; \textbf{tidak dapat dilewati}; \textbf{wajib} untuk semua form publik.
	\end{block}
	Gunakan keduanya: client-side untuk UX, server-side untuk keamanan.
\end{frame}

% ========== SECTION 6-5: Elemen Semantik ==========
\begin{frame}{Struktur Halaman dengan Elemen Semantik}
	\begin{itemize}
		\item \texttt{<header>}: kepala halaman/section; \texttt{<nav>}: navigasi utama.
		\item \texttt{<main>}: konten utama (\textbf{satu per halaman}).
		\item \texttt{<article>}: konten mandiri; \texttt{<section>}: grup tematik; \texttt{<aside>}: konten pendukung.
		\item \texttt{<footer>}: penutup halaman.
	\end{itemize}
	\begin{block}{Manfaat}
		Makna semantik $\to$ SEO lebih baik, pembaca layar navigasi otomatis, kode lebih mudah dirawat.
	\end{block}
\end{frame}

% ========== SECTION 6-6: Media ==========
\begin{frame}{Media: Audio dan Video}
	\begin{itemize}
		\item \texttt{<audio>}, \texttt{<video>}: native HTML5 tanpa plugin.
		\item Atribut: \texttt{controls}, \texttt{autoplay}, \texttt{loop}, \texttt{muted}; \texttt{poster} (video thumbnail).
		\item \texttt{<source>}: beberapa format fallback (MP3/WAV, MP4/WebM).
		\item \texttt{<track>}: subtitle (\texttt{kind="subtitles"}), caption (\texttt{kind="captions"}); format WebVTT (\texttt{.vtt}).
	\end{itemize}
	\begin{block}{Aksesibilitas media}
		Teks alternatif, transkrip, subtitle untuk pengguna tuli/tuna rungu.
	\end{block}
\end{frame}

% ========== SECTION 6-7: SVG ==========
\begin{frame}{Gambar Vektor (SVG) dalam HTML}
	\begin{itemize}
		\item \textbf{SVG} vs raster: vektor skalabel tanpa kehilangan kualitas.
		\item SVG inline: sisipkan \texttt{<svg>} langsung di HTML.
		\item Aksesibilitas: \texttt{<title>}, \texttt{<desc>}, \texttt{aria-label}, \texttt{role="img"}.
	\end{itemize}
\end{frame}

% ========== SECTION 6-8: Embedding ==========
\begin{frame}{Embedding: iframe, object, dan embed}
	\begin{itemize}
		\item \texttt{<iframe>}: menyematkan halaman/konten eksternal.
		\item \texttt{sandbox}: membatasi aksi (script, form, dll.) untuk keamanan.
		\item \texttt{allow}: izinkan fitur (camera, geolocation, dll.).
		\item \texttt{title}: deskripsi untuk aksesibilitas.
	\end{itemize}
	\begin{block}{Keamanan}
		Gunakan \texttt{sandbox} dengan pembatasan yang tepat; jangan embed sumber tidak tepercaya.
	\end{block}
\end{frame}

% ========== SECTION 6-9: Aksesibilitas Prinsip ==========
\begin{frame}{Aksesibilitas (a11y): Prinsip dan Praktik}
	\begin{itemize}
		\item Web dapat digunakan \textbf{semua orang} (disabilitas visual, auditori, motorik, kognitif).
		\item Standar \textbf{WCAG} (W3C); prinsip \textbf{POUR}:
		\item \textbf{Perceivable}: konten dapat dipersepsi (teks alt, kontras warna).
		\item \textbf{Operable}: antarmuka dapat dioperasikan (keyboard, fokus).
		\item \textbf{Understandable}: informasi dapat dipahami (bahasa, error jelas).
		\item \textbf{Robust}: kompatibel berbagai teknologi (HTML semantik, ARIA).
	\end{itemize}
	Level: A (minimum), AA (direkomendasikan), AAA (optimal).
\end{frame}

\begin{frame}{ARIA: Pelengkap Semantik HTML}
	\begin{itemize}
		\item \textbf{ARIA} (Accessible Rich Internet Applications): bila HTML semantik tidak cukup.
		\item \texttt{role}: peran elemen (\texttt{role="dialog"}), \texttt{aria-label}, \texttt{aria-describedby}.
		\item \texttt{aria-expanded}, \texttt{aria-hidden}, \texttt{aria-live}, dll.
	\end{itemize}
	\begin{block}{Aturan pertama ARIA}
		Jangan gunakan ARIA jika elemen HTML semantik sudah cukup (\texttt{<button>}, \texttt{<nav>}, dll.).
	\end{block}
\end{frame}

% ========== SECTION 6-10: Aksesibilitas Tabel dan Form ==========
\begin{frame}{Aksesibilitas Tabel dan Form}
	\begin{itemize}
		\item \textbf{Tabel}: \texttt{<caption>}, \texttt{scope}, \texttt{headers}; hindari sel kosong tanpa konteks.
		\item \textbf{Form}: \texttt{<label for="...">}, \texttt{<fieldset>}/\texttt{<legend>}, \texttt{aria-describedby} untuk pesan error.
		\item Semua elemen interaktif dapat diakses via \textbf{keyboard}; urutan tab logis.
		\item Uji dengan pembaca layar (NVDA, VoiceOver) atau navigasi keyboard saja.
	\end{itemize}
\end{frame}

% ========== SECTION 6-11: Validasi HTML ==========
\begin{frame}{Validasi HTML dan Alat Validator}
	\begin{itemize}
		\item \textbf{W3C Markup Validation Service}: \texttt{https://validator.w3.org/}
		\item \textbf{Error}: harus diperbaiki (struktur tidak valid).
		\item \textbf{Warning}: perbaikan disarankan (deprecated, best practice).
		\item Dokumen valid $\to$ interoperabilitas, aksesibilitas, perawatan lebih mudah.
	\end{itemize}
\end{frame}

% ========== SECTION 6-12: Debugging ==========
\begin{frame}{Debugging HTML: Menemukan dan Memperbaiki Error}
	\begin{itemize}
		\item \textbf{DevTools} (F12): tab \textbf{Elements} (inspeksi HTML/CSS), \textbf{Console} (error JS).
		\item Periksa struktur DOM, atribut hilang, tag tidak tertutup.
		\item Validasi W3C untuk error markup; Lighthouse untuk skor aksesibilitas.
	\end{itemize}
\end{frame}

% ========== SECTION 6-13: Best Practices ==========
\begin{frame}{Best Practices HTML dan Standar (WHATWG / W3C)}
	\begin{itemize}
		\item Ikuti standar \textbf{WHATWG} (HTML Living Standard) dan \textbf{W3C}.
		\item Dokumen valid, elemen semantik tepat, atribut aksesibilitas lengkap.
		\item Interoperabilitas lintas browser; kode profesional dan dapat dirawat.
	\end{itemize}
\end{frame}

% ========== Rangkuman dan Penutup ==========
\begin{frame}{Rangkuman Bab 6}
	\begin{itemize}
		\item Tabel: \texttt{thead}/\texttt{tbody}/\texttt{tfoot}, \texttt{scope}, \texttt{caption}; aksesibilitas.
		\item Form: \texttt{form}, \texttt{input}, \texttt{label}, \texttt{button}; tipe HTML5; validasi client + server; \texttt{fieldset}/\texttt{legend}.
		\item Elemen semantik: \texttt{header}, \texttt{nav}, \texttt{main}, \texttt{article}, \texttt{section}, \texttt{aside}, \texttt{footer}.
		\item Media: \texttt{audio}/\texttt{video}, \texttt{<track>} subtitle; SVG inline; iframe dengan sandbox.
		\item Aksesibilitas: POUR (WCAG), ARIA bila perlu; label, caption, keyboard.
		\item Validasi W3C; debugging DevTools; best practices WHATWG/W3C.
	\end{itemize}
\end{frame}

\begin{frame}{}
	\centering
	\vfill
	\textbf{\Large Pertanyaan?}
	\vfill
\end{frame}

\end{document}
